\section{Circle collapse and integral homology}\label{sec:circle-collapse}

We first compute the inner page in the \(h\)-model and prove acyclicity
of the large total. These facts allow a complete comparison with an
ordinary local Milnor tube. We then cut the exterior page in the cone
model to determine which classes in \(H\otimes E\) extend across
the inner piece. Section~\ref{sec:integral-lattice} will determine
the inner coordinate of each integral extension.

\subsection{The circle collapse}

\begin{lemma}[inner homology and acyclic total]\label{lem:acyclic}
There are natural identifications
\begin{equation}
H_3(W)=K,\qquad H_2(W)=H/\pi_*D,\qquad
\Psi:H_3(W;\Z/k)\hookrightarrow D/kD.
\label{eq:tube-11}
\end{equation}
The large total \(\widehat\T_s\) is acyclic.
\end{lemma}
\begin{proof}
Replace \(h+P\) by \(h\) on the inner piece using
Lemma~\ref{lem:continuations}. The base-with-\(w\) page is
\[
G=(\mathcal B\cap\{h=c\})\times\{|w|\le\lambda\},
\]
and its root subpage \(\D\) is a page of \(d\). For \(R_s\ne0\),
we retract the vertical hyperbola to a circle; over the root page,
that circle collapses. We realize this model by balancing the two
moduli inside the inner block. For fixed \((p,w)\), write \(A=|y|^2/\sigma^2\),
\(B=\sigma^2|z|^2\), so \(AB=|R_s|^2\), and choose
\(\sigma\) near \(\lambda^{(\alpha-\beta)/2}\).
For \(0\le t\le1\), replace \(A-B\) by
\((1-t)(A-B)\) while retaining \(AB\) and the circle angle.
The positive solution for \(A,B\) is continuous at \(R_s=0\);
on an axis it shrinks the active coordinate to zero.
Both moduli stay between their starting and final values.
Since \(\max|R_s|\le(1+o(1))\lambda^k\), the core
\(|y|=\sigma\sqrt{|R_s|}\),
\(|z|=\sigma^{-1}\sqrt{|R_s|}\) stays strictly within the
\(M\)-bounds. This is a deformation retraction, relative to
the constant-angle section, onto
\begin{equation}
X(G,\D)=(G\times S^1)/
(\{g\}\times S^1\text{ collapsed for each }g\in\D).
\label{eq:tube-12}
\end{equation}
The constant-angle section splits off the homology of \(G\).
Quotienting that section gives \((G/\D)\wedge S^1\), so the
remaining homology is \(H_{*-1}(G,\D)\), naturally for moving roots.
The spaces \(G\) and \(\D\) are connected and have only first
reduced homology. Their relative long exact sequence therefore gives
\[
0\longrightarrow H_2(G,\D)\longrightarrow D
\xrightarrow{\pi_*}H\longrightarrow H_1(G,\D)\longrightarrow0.
\]
Thus the underlying groups are \(H_3(W)\cong K\) and
\(H_2(W)\cong H/\pi_*D\). We next fix the signs of these
identifications and of the mod-\(k\) injection; these signs determine
the residue condition in Theorem~\ref{thm:tube}.

Orient this \(S^1\) by increasing \(\arg y\).
For an arc \(I\) from \(w_0\) to a root, use the circle-first
orientation \([S^1_y]\times I\) on its sweep; the collapsed root
end leaves boundary \(+\gamma\) at \(w_0\).
Write \(\mathfrak S(C)=[S^1_y]\times C\) for circle-first suspension.
The cross-product boundary rule gives
\[
 \partial\mathfrak S(C)=-\mathfrak S(\partial C),\qquad
 \mathfrak S(a_0\times I)=-a_0\times([S^1_y]\times I)
\]
in a product chart over an oriented base interval \(a_0\), with
its endpoint faces taken relatively. The root part of the
boundary of \(a_0\times I\) is \(-[a_{0,\rm root}]\); write
\(\partial_{\D}\) for this root-endpoint boundary, and for the
resulting relative connecting map on relative cycles.
Define \(\Psi\) as the \emph{negative} of the relative root-endpoint
connecting map under this circle-first suspension, namely
\(-\partial_{\D}\mathfrak S^{-1}\). Thus the local chain
\(-a_0\times(S^1_y\times I)\) has signed root contribution
\(+[a_{0,\rm root}]\). The interval endpoint faces and the
remaining \(w_0\)-boundary must cancel before a homology class is
defined. The next section supplies this cancellation by a lens
chain with a cap at the seam, whose local root vector is diagonal.
Then \(\Psi\) applies to the resulting mod-\(k\) class with the
sign just computed; no individual root is assumed to close around
the whole base loop.
This fixes the sign of \(\Psi\) throughout.
We use the same suspension convention for \(H_3(W)=K\).
In degree two, identify \(H_2(W)\) with \(H/\pi_*D\) by the
negative of circle-first unsuspension. Künneth products on
\(J\simeq\Hh\times L\) are ordered with the base factor first:
for a closed oriented base loop \(a\),
\([a\times\gamma]=-\mathfrak S(a)\). This normalization therefore
sends the interface torus \([a\times\gamma]\) to
\(+a\bmod\pi_*D\). The quotient is defined on \(H/kH\) because
\(kH\subseteq\pi_*D\), by \(\pi_*\pi^!=k\).
With \(\Z/k\) coefficients, the same relative sequence injects
\(H_2(G,\D;\Z/k)\) into \(H_1(\D;\Z/k)=D/kD\).
Under suspension, the signed injection from \(H_3(W;\Z/k)\)
is the negative relative boundary just fixed, namely \(\Psi\)
in \eqref{eq:tube-11}.

It remains to prove acyclicity of the large total. The same collapse
argument for total disks first shows that \(\widehat W\) is acyclic,
because the total tubes of \(h,d\) are acyclic. We control the
exterior attachment using the explicit contraction of the cone model.

At \(s=0\), put \(d_P=1+a\alpha=1+b\beta\).
The explicit contraction
\begin{equation}
(X,T;y,z,w)\longmapsto
\bigl(t^{2/5}X,t^{1/5}T;\
t^{\alpha/d_P}y,t^{\beta/d_P}z,t^{1/d_P}w\bigr),
\qquad 0\le t\le1,
\label{eq:tube-12a}
\end{equation}
scales both \(h\) and \(P\) by \(t\), and scales
both \(yz\) and \(w^k\) by \(t^{k/d_P}\).
It preserves the base ball, its \(|h|\)-cut,
the vertical block cuts, and the \(|h+P|\)-cut.
Hence \(\widehat\T_0\) and its inner part are
contractible. Mayer--Vietoris therefore makes
\(\widehat J_0\to\widehat M_{{\rm ext},0}\) a homology
equivalence. Transport by Lemma~\ref{lem:continuations} gives
the same statement at positive \(s\). Attaching the acyclic
\(\widehat W\) along \(\widehat J\) leaves
\(\widehat\T_s\) acyclic.
\end{proof}

\subsection{Comparison with an ordinary Milnor tube}

Acyclicity lets us identify page homology through relative homology:
for an acyclic total, the connecting map identifies relative \(H_4\)
with page \(H_3\). We compare the large and ordinary total/page pairs
first, retaining the filling chains that determine the ordered form.

\begin{lemma}[The large tube computes the local page]\label{lem:local}
The middle homology and filling pairing of the page of \eqref{eq:tube-4}
are naturally those of a sufficiently small ordinary Milnor
representative of \(f_s\).
\end{lemma}
\begin{proof}
Fix \(s>0\) after Lemmas~\ref{lem:boundary} and
\ref{lem:isolation}. In a holomorphic chart on the smooth germ
\(C_s\), choose a closed ordinary Milnor ball \(B_{\rm m}\) inside
the interior of every outer cut of \(X_s\). Choose
\(0<\delta'\le\delta\) so that its sphere is transverse to all
fibers over \(\{|v|\le\delta'\}\). The corresponding ordinary
analytic total
\[
 \widehat N_s=B_{\rm m}\cap f_s^{-1}(\{|v|\le\delta'\})
\]
is contractible; this is the ordinary tube consequence of Milnor's
proof recorded in \cite[Theorem~1.8 and Remark~1.9]{Massey2015LeCycles}.
It does not apply by itself to the custom outer cuts.

Let \(\widehat\T_s'\) be the large total with value radius
\(\delta'\), and set
\[
 \widehat A_s=(X_s\setminus\operatorname{int}B_{\rm m})
              \cap f_s^{-1}(\{|v|\le\delta'\}),\qquad
 \widehat S_s=\partial B_{\rm m}
              \cap f_s^{-1}(\{|v|\le\delta'\}).
\]
Write \(N_s,A_s,S_s,\T_s'\) for the corresponding pages at a fixed
\(\lvert c'\rvert=\delta'\). Stratify \(\widehat A_s\) by the
genuine outer faces and the new Milnor-sphere face. The auxiliary
cut \(J\) remains ordinary interior in this comparison.
Lemma~\ref{lem:boundary} gives rank two on the outer strata over
the whole disk, the Milnor choice gives it on the sphere, and
Lemma~\ref{lem:isolation} gives it in the intervening interior.
The cuts are compact, so \(f_s\) is a proper stratified submersion
on \(\widehat A_s\) and on \(\widehat S_s\) over this disk,
stratified by its interior and boundary circle.
The face-preserving Ehresmann argument over the contractible disk
therefore makes \(A_s\hookrightarrow\widehat A_s\) and
\(S_s\hookrightarrow\widehat S_s\) homotopy equivalences.
Their relative homology groups vanish. Relative Mayer--Vietoris
for
\[
 (\widehat\T_s',\T_s')
   =(\widehat N_s,N_s)\ \cup\ (\widehat A_s,A_s),
 \qquad
 (\widehat N_s,N_s)\cap(\widehat A_s,A_s)
   =(\widehat S_s,S_s)
\]
now identifies the large total/page relative homology with the
ordinary local one.

The disk-wide outer regularity and isolation away from value zero
also permit the large radius to shrink from \(\delta\) to
\(\delta'\) through all positive intermediate radii: horizontal
collar fields are tangent to every outer face and vanish near the
critical region. Lemma~\ref{lem:acyclic} proves that
\(\widehat\T_s\) is acyclic at the original radius; this isotopy
therefore makes \(\widehat\T_s'\) acyclic. The ordinary total
\(\widehat N_s\) is contractible as above. The connecting maps in
the relative long exact sequences of these two acyclic totals identify
their relative \(H_4\)-groups with the page \(H_3\)-groups, proving
the homology assertion. A local filling can be used unchanged
in the large total. Transporting the two page cycles from the
smaller to the original value circle along two disjoint angular
rays adds collars at disjoint complex values, hence no new
intersections. Thus the filling pairing is identical.
\end{proof}

\subsection{The signed exterior cut}

For the gluing, the inner page has middle group \(K\), and a boundary torus
\([a\times\gamma]\) dies there precisely when \(a\in\pi_*D\).
We next compute the boundary torus carried by an exterior cycle.

The exterior page admits a useful two-sided cut. We normalize
the value \(c\) to lie on the positive real axis; the picture
shows the \(q=h\) value disk, not the spatial cone block.

\begin{figure}[ht]
\centering
\begin{tikzpicture}[x=1cm,y=1cm,>=Latex]
  \draw[->] (-4.15,1.5)--(4.3,1.5)
    node[right] {\(\operatorname{Re}(q/c)\)};
  \foreach \x/\lab in {-2/0,0/{1/2},2/1}
    {\draw (\x,1.4)--(\x,1.6) node[above] {\(\lab\)};}
  \fill[blue!10] (-4,0) rectangle (0,1);
  \fill[orange!16] (0,0) rectangle (4,1);
  \draw (-4,0) rectangle (4,1);
  \draw[dashed,thick] (0,-.07)--(0,1.15);
  \node[align=center] at (-2,.53)
    {\(P=c-q\ne0\)\\ side retracts to \(Q\)};
  \node[align=center] at (2,.53)
    {radial contraction\\ stops at \(\partial Y_\lambda\)};
  \node[below] at (0,-.12) {overlap \(\simeq\Hh\times Q\)};
  \draw[->,thick] (2,-.07)--(2,-.7);
  \node[below,align=center] at (2,-.75)
    {\(J\simeq\Hh\times L\)\\ attaches to inner \(W\)};
\end{tikzpicture}

\caption{The exterior cut uses \(q=h\), with the full page equation
\(P=c-q\), shown by its real-part projection. The left side keeps
\(P\) away from zero; the right
side retracts to the actual interface \(J\) by stopping weighted
radial contraction at \(\partial Y_\lambda\), where \(P\) need not
vanish. The interface continuation identifies \(J\) with
\(\Hh\times L\). The overlap-to-\(J\) map is
\(\mathrm{id}_{\Hh}\times(Q\to L)\). This is schematic:
the disk radius \(\varepsilon\) and point \(c\) are not to scale.}
\label{fig:split}
\end{figure}

\begin{lemma}[Mayer--Vietoris boundary]\label{lem:mv}
The full page has an exact sequence
\begin{equation}
0\longrightarrow K\longrightarrow H_3(\T_s)
\xrightarrow{\partial} H\otimes E
\xrightarrow{\overline{1\otimes\tau}}H/\pi_*D.
\label{eq:tube-13}
\end{equation}
The image of the exterior page in \(H_3(\T_s)\) maps
isomorphically under \(\partial\) to \(H\otimes\ker\tau\).
\end{lemma}
\begin{proof}
Use Lemma~\ref{lem:continuations} to compute the exterior at \(s=0\).
Normalize \(c>0\), and put \(q=h\). Cut at
\(\operatorname{Re}q=c/2\), as in Figure~\ref{fig:split}, calling
the left and right pieces \(A\) and \(B\), respectively.
On \(A\), \(|P|=|c-q|\ge c/2\), so the removed inner block is
absent. The allowed \(P\)-values form a convex region avoiding zero,
and the boundary estimates of Lemma~\ref{lem:boundary} trivialize
its \(P\)-pages with their outer cuts. Contracting the base by the
positive weights of \(h\) and carrying these pages through the
trivialization retracts \(A\) to \(Q\).

On \(B\), choose one boundary-respecting trivialization of the
\(h\)-pages over the convex regular-value region
\(\{|q|\le\varepsilon,\ \operatorname{Re}q\ge c/2\}\).
It can agree near \(c\) with the page labels used in
Lemma~\ref{lem:continuations}. For \(v=(y,z,w)\), define the
radius of its block by
\[
 r(v)=\max\left\{
 (|y|/M)^{1/\alpha},\ (|z|/M)^{1/\beta},\ |w|
 \right\}.
\]
The exterior condition is \(\lambda\le r(v)\le\rho\).
Scale \(v\) by the cone weights, from \(t=1\) down to
\(t=\lambda/r(v)\), while retaining its \(h\)-page label and setting
\[
 v_t=(t^\alpha y,t^\beta z,tw),\qquad
 q_t=c-t^{d_P}P(v)=(1-t^{d_P})c+t^{d_P}q,
 \qquad d_P=1+a\alpha.
\]
This keeps \(\lambda\le r(v_t)\le\rho\), and \(q_t\) stays in
the same convex \(h\)-region. The base-ball cut is preserved by
the trivialization. The homotopy therefore stays in \(B\) and
stops at the actual interface \(J\), fixing \(J\) pointwise.
At its endpoint \(v_L\in L\), the value is
\(q=c-P(v_L)\), generally different from \(c\).
The \(s=0\) interface continuation of
Lemma~\ref{lem:continuations} retains \(v_L\) and the base label
while changing this value to \(c\). It identifies the retract
with \(\Hh\times L\).

The seam values form the contractible interval
\(I=\{q=c/2+iu:|q|\le\varepsilon\}\).
Trivializing both page factors along \(I\) identifies the overlap
\(C=A\cap B\) with \(\Hh\times Q\times I\).
After contracting \(I\), its map to \(J\simeq\Hh\times L\)
is homotopic to \(\mathrm{id}_{\Hh}\times(Q\to L)\), where
\(Q\to L\) is weighted radial projection.
Both \(\Hh,Q\) have homotopy dimension one, so overlap
\(H_3=0\), while overlap \(H_2=H\otimes E\).
Moreover \(H_2(J)=H\otimes\Z/k=H/kH\) and
\(H_3(J)=\Z\). The latter generator maps to zero in
\(H_3(W)=K\): it is the full inner fiber over a regular
base point, whose \(H_3\) vanishes in the collapse model.
We record the orientation of the remaining interface map.
For a full-page
3-cycle write its exterior chain as \(e=e_A+e_B\), with
the inner chain \(w\) satisfying
\(\partial w=b_J,\ \partial e=-b_J\).
Define \(t_C=-\partial e_A\in Z_2(C)\). Then
\begin{equation}
\partial e_B=t_C-b_J,\qquad
z:=[t_C]\in H_2(C)=H\otimes E.
\label{eq:mv-cut-chain}
\end{equation}
This fixes the sign of the Mayer--Vietoris coordinate.
Let \(r_t:C\to B\), \(0\le t\le1\), be the stopped radial
homotopy just constructed, reparametrized to end in \(J\).
Its endpoint is the actual interface, with the same \(h\)-page
label under the compatible parameter rectangle. For the
cycle \(t_C\), the interval-first prism chain
\(S=r_*([0,1]\times t_C)\) has
\(\partial S=r_{1*}t_C-t_C\). Consequently
\begin{equation}
\partial(e_B+S)=r_{1*}t_C-b_J,\qquad
[b_J]=r_{1*}[t_C]
     =(1\otimes\tau)z\quad\hbox{in }H_2(J)=H/kH.
\label{eq:mv-strip}
\end{equation}
Under the degree-two normalization fixed after
\eqref{eq:tube-12}, this is the positive sign of the last arrow
in \eqref{eq:tube-13}.
The Mayer--Vietoris sequence now gives \eqref{eq:tube-13}, with final
map the quotient of \(1\otimes\tau\). Before \(W\) is
attached, exterior \(H_3\) differs from its
\(H\otimes\ker\tau\) contribution by precisely the image
of \(H_3(J)\); that image dies after attachment. This
proves the last statement.
\end{proof}
